% !TeX root = ../navier-stokes-manuscript.tex
\subsection{Where viscosity pays on a strict interval}\label{sub:SobolevRegularityInsideEveryCompletedRow}

Fix one response \(V_r\) and one spatial height \(s\).  Its scalar energy and
complete transfer are
\[
  E_{r,s}:=\frac12\norm{\Lambda^sV_r}_2^2
\]
and
\[
  \mathcal P_{r,s}
  :=-
  \operatorname{Re}
  \pair{\Lambda^sV_r}
  {\Lambda^s\!\left[
    \sum_{a=0}^{r}\binom{r}{a}(V_a\cdot\nabla)V_{r-a}
    +\nabla Q_r
  \right]}_{L^2}.
\]
Pairing the \(r\)-th response equation at this height gives
\begin{equation}\label{eq:BodyEveryTemporalRowSpatialEnergy}
  E_{r,s}'=\mathcal P_{r,s}-2\nu E_{r,s+1}.
\end{equation}
Beginning at \(s=\tfrac12\) and repeating the induction in
\Cref{prop:CriticalHeightSpatialAscent} produces
\begin{equation}\label{eq:BodyEveryTemporalRowCompletedHeight}
  \partial_t^nE_{r,1/2}
  =(-2\nu)^nE_{r,n+1/2}
  +\sum_{j=0}^{n-1}
  (-2\nu)^j
  \partial_t^{\,n-1-j}\mathcal P_{r,j+1/2}.
\end{equation}
The pattern repeats inside every response, together with the same obstruction:
the higher spatial energy appears, and the nonlinear transfer remains beside
it.

To see an actual spatial gain, stop first at one strict cutoff
\(T<T_*\) and write \(I=(0,T)\).  Apply the Leray projection to the
differentiated momentum equation.  It removes the pressure gradient from this
estimate while leaving the pressure-complete recurrence intact:
\[
  \partial_tV_r-\nu\Delta V_r=-F_r,
  \qquad
  F_r:=\Leray\sum_{a=0}^{r}\binom{r}{a}(V_a\cdot\nabla)V_{r-a}.
\]
For an integer \(m\geq2\), the finite-order product bound proved in
\Cref{prop:LocalMaximalHistory} gives
\begin{equation}\label{eq:BodyTemporalResponseForcingEstimate}
  \norm{F_r}_{H^{m-1}}
  \leq
  C_{m,r}\sum_{a=0}^{r}
  \norm{V_a}_{H^m}\norm{V_{r-a}}_{H^m}.
\end{equation}
The same local construction preserves every
finite derivative already present in the smooth datum on this chosen interval.
For the finite indices now in use,
\begin{equation}\label{eq:BodyStrictIntervalResponseBound}
  \max_{0\leq a\leq r}
  \norm{V_a}_{L^\infty(0,T;H^m)}<\infty.
\end{equation}
Hence the forcing in
\eqref{eq:BodyTemporalResponseForcingEstimate} is square-integrable in time.

Pair the projected equation one spatial step above its forcing.  The heat term
now makes the payment that the triangular identity only exposed:
\begin{equation}\label{eq:BodyTemporalResponseParabolicGain}
  \frac d{dt}\norm{V_r}_{H^m}^2
  +\nu\norm{V_r}_{H^{m+1}}^2
  \leq
  C_{\nu,m}\left(
    \norm{V_r}_{H^m}^2+\norm{F_r}_{H^{m-1}}^2
  \right).
\end{equation}
After time integration, \(V_r\) gains one spatial degree.  Returning to the
equation puts its next time response one degree below:
\begin{equation}\label{eq:BodyStrictIntervalAdjacentResponsePair}
  V_r\in L^2(I;H^{m+1}),
  \qquad
  V_{r+1}=\partial_tV_r\in L^2(I;H^{m-1}).
\end{equation}
This is the first place where the spatial movement has been estimated rather
than read from an identity.

The finite estimates may be placed beside one another as a two-dimensional
grid.  Moving down changes the temporal response; moving across changes the
spatial resolution of that response.
\begin{center}
\captionsetup{hypcap=false}
\captionof{table}{The mixed derivative grid.  Each dot names a coordinate of the same velocity; no dot is a cutoff-uniform endpoint estimate.}
\label{tab:MixedDerivativeGrid}
\renewcommand{\arraystretch}{1.35}
\begin{tabular}{@{}c c c c c c@{}}
\toprule
& \(H^0_x\) & \(H^1_x\) & \(H^2_x\) & \(H^3_x\) & \(\cdots\) \\
\midrule
\(V_0=u\) & \(\bullet\) & \(\bullet\) & \(\bullet\) & \(\bullet\) & \(\cdots\) \\
\(V_1=u_t\) & \(\bullet\) & \(\bullet\) & \(\bullet\) & \(\bullet\) & \(\cdots\) \\
\(V_2=u_{tt}\) & \(\bullet\) & \(\bullet\) & \(\bullet\) & \(\bullet\) & \(\cdots\) \\
\(V_3=u_{ttt}\) & \(\bullet\) & \(\bullet\) & \(\bullet\) & \(\bullet\) & \(\cdots\) \\
\(\vdots\) & \(\vdots\) & \(\vdots\) & \(\vdots\) & \(\vdots\) & \(\ddots\) \\
\bottomrule
\end{tabular}
\end{center}

Now ask for a finite observation of the solution: \(q\) time derivatives and
\(k\) spatial derivatives.  Choose finite heights
\[
  R>q+\frac12,
  \qquad
  M>k+\frac32.
\]
The vector-valued time embedding and the three-dimensional spatial embedding
\parencite[Theorem~7.28]{Cerda2010} give
\[
  H^R\bigl(I;H^M(\R^3)\bigr)
  \hookrightarrow
  C^q\bigl(\overline I;C^k(\R^3)\bigr).
\]
Repeating the finite payment high enough in each direction fills every finite
request on this fixed \(I\).  Collecting those requests on the same history
gives
\begin{equation}\label{eq:BodyFormalMixedIntersection}
  \mathscr V_\infty(I)
  :=\bigcap_{r=0}^{\infty}\bigcap_{m=0}^{\infty}
  H^r\bigl(I;H^m(\R^3)\bigr)
  \hookrightarrow C^\infty(\overline I\times\R^3).
\end{equation}
The constants may depend on the chosen finite coordinate.  Smoothness asks for
each such coordinate in turn, not for one constant uniform over the whole
grid.

The cutoff is the unpaid part.  Every dot is finite after \(T<T_*\) has been
fixed, but none of the estimates above prevents its constant from diverging as
\(T\uparrow T_*\).  To carry one spatial coordinate to that time slice, it
must be paired with the temporal response directly below it.  Those adjacent
rows form the finite rectangles.
